% !TEX root = ../main.tex
% Resultados literales de la ejecucion del 06/08/2026.
% Fuente: anexo del ADR-0003, escrito por el propio script.
% NO editar estas cifras a mano: si cambian, se vuelve a ejecutar
%   py scripts/experimentos/experimento_embeddings.py
% y se copian de ahi. Corpus de 1.334 fragmentos y 50 preguntas.
\begin{table}[htbp]
  \centering
  \small
  \caption[Comparativa de modelos de incrustaciones]{Comparativa de modelos de
    incrustaciones sobre la colección completa y las 50 preguntas del
    conjunto de evaluación.
    Fuente: elaboración propia.}
  \label{tab:comparativa_embeddings}
  % Ocho columnas y un nombre de modelo largo en monoespaciada desbordaban la
  % caja 14,6 pt. Se aprieta la separacion entre columnas en lugar de encoger
  % la letra: 2 pt menos a cada lado de ocho columnas son 32 pt, de sobra.
  \setlength{\tabcolsep}{4pt}
  \begin{tabular}{@{}lrrrrrrr@{}}
    \toprule
    \textbf{Modelo}                            & \textbf{R@3}   & \textbf{R@5}   & \textbf{R@10}
                                               & \textbf{RU@3}  & \textbf{RU@5}  & \textbf{RU@10} & \textbf{MRR}   \\
    \midrule
    \texttt{multilingual-e5-small}             & 0,697          & 0,803          & 0,911
                                               & 0,985          & 0,990          & \textbf{1,000} & \textbf{0,970} \\
    \texttt{multilingual-e5-large}             & \textbf{0,740} & \textbf{0,852} & \textbf{0,938}
                                               & \textbf{0,995} & \textbf{1,000} & \textbf{1,000} & \textbf{0,970} \\
    \texttt{bge-m3}                            & 0,720          & 0,836          & 0,917
                                               & 0,985          & 0,995          & 0,995          & 0,949          \\
    \texttt{sentence\_similarity\_spanish\_es} & 0,152          & 0,194          & 0,307
                                               & 0,410          & 0,505          & 0,610          & 0,342          \\
    \bottomrule
  \end{tabular}
  \begin{minipage}{0.95\textwidth}
    \vspace{4pt}
    \tiny
    \textbf{Notas:} \textbf{R@K} es la exhaustividad por \emph{fragmento} (cuántos de los
    trozos de la unidad correcta aparecen entre los $K$ primeros resultados).
    \textbf{RU@K} es la exhaustividad por \emph{unidad} (si se ha encontrado la asignatura
    correcta, sin penalizar que falte alguno de sus trozos). \textbf{MRR}
    (\textit{Mean Reciprocal Rank}) es la media del inverso de la posición del primer resultado
    correcto, de modo que premia colocarlo antes. En negrita, el mejor valor de cada
    columna. El
    modelo adoptado es \texttt{multilingual-e5-small}.
  \end{minipage}
\end{table}
